% ============================================================
%  Chương VIII – Kết luận
% ============================================================
\chapter{Kết luận}
\label{chap:conclusion}

% ----------------------------------------------------------
\section{Tổng kết}
\label{sec:8.1}

Báo cáo đã đi từ nền tảng mật mã khóa công khai ở
Chương~\ref{chap:crypto} đến bài toán phân phối khóa và ràng buộc danh tính
trong hạ tầng PKI ở Chương~\ref{chap:pki-arch}. Trên cơ sở đó,
Chương~\ref{chap:x509} trình bày chứng chỉ số, X.509 và các chuẩn liên quan,
trước khi Chương~\ref{chap:deploy} phân tích cách PKI xuất hiện trong HTTPS,
chữ ký số, S/MIME, VPN và code signing.

Các chương cuối đặt PKI vào bối cảnh vận hành. Chương~\ref{chap:opensource}
khảo sát các công cụ và nền tảng mã nguồn mở, từ OpenSSL, XCA đến EJBCA,
Certbot, Vault PKI và FreeIPA. Chương~\ref{chap:challenges} tổng hợp các
thách thức về vòng đời chứng chỉ, thu hồi, niềm tin vào CA, tự động hóa và
crypto-agility. Kết luận chính là PKI không chỉ gồm thuật toán hay
certificate. Nó là hệ thống kết hợp mật mã, chính sách tin cậy, quy trình vận
hành và công cụ quản trị.

% ----------------------------------------------------------
\section{Ý nghĩa và lời cảm ơn}
\label{sec:8.2}

Qua quá trình thực hiện báo cáo, nhóm hiểu rõ hơn cách public key được gắn với
danh tính thông qua CA và certificate chain. Báo cáo cũng giúp nhóm nhìn thấy
vai trò của revocation, trust store, certificate lifecycle và tự động hóa trong
một hệ PKI thực tế. Những nội dung này tạo nền tảng để tiếp tục đọc tài liệu
kỹ thuật hoặc xây dựng mô hình thử nghiệm PKI ở mức có kiểm soát.

Nhóm xin chân thành cảm ơn PGS. TS. Nguyễn Linh Giang đã hướng dẫn, cung cấp
định hướng môn học và hỗ trợ nhóm trong quá trình thực hiện báo cáo.
